% BHU PYQ -- Geography: Environmental Ethics and Sustainability (2025-26 NEP)
\documentclass[11pt,a4paper]{article}
\usepackage[a4paper,top=2.5cm,bottom=2.5cm,left=2.8cm,right=2.8cm,headheight=14pt]{geometry}
\usepackage{amsmath,amssymb}
\usepackage{fontspec}
\setmainfont{Kohinoor Devanagari}
\usepackage{microtype,fancyhdr,enumitem,hyperref,lastpage,parskip}

\hypersetup{colorlinks=true,urlcolor=black,linkcolor=black,pdftitle={GGRMD21: Environmental Ethics and Sustainability -- BHU NEP 2025-26}}

\pagestyle{fancy}\fancyhf{}
\renewcommand{\headrulewidth}{0.4pt}\renewcommand{\footrulewidth}{0.4pt}
\fancyhead[L]{\small \textit{Banaras Hindu University}}
\fancyhead[R]{\small \textit{GGRMD21: Environmental Ethics and Sustainability (2025-26)}}
\fancyfoot[C]{\small Page \thepage\ of \pageref{LastPage}}

\newlist{parts}{enumerate}{2}
\setlist[parts,1]{label=(\alph*),leftmargin=*,itemsep=4pt,topsep=3pt}
\setlist[parts,2]{label=(\roman*),leftmargin=*,itemsep=2pt,topsep=2pt}

\newcommand{\pts}[1]{\hfill \textit{[#1]}}

\begin{document}

\begin{center}
  {\large\bfseries BANARAS HINDU UNIVERSITY}\\[2pt]
  {\normalsize Under Graduate Semester II Examination 2025-26 (NEP)}\\[6pt]
  \rule{0.8\linewidth}{0.5pt}\\[6pt]
  {\large\bfseries GEOGRAPHY}\\[4pt]
  {\large\bfseries Paper Code: GGRMD21 : Environmental Ethics and Sustainability}\\[4pt]
  {\normalsize Time: 2:30 Hours \hfill Max. Marks: 70}
\end{center}
\rule{\linewidth}{0.4pt}
\smallskip

\noindent \textit{Instructions: Attempt five questions including Question 1 which is compulsory. Answer each part of Question 1 in approximately 200 words and of remaining questions in approximately 1000 words. The figures in the right-hand margin indicate marks. Illustrate your answers with suitable maps and diagrams. Write your Roll No.\ at the top immediately on receipt of this question paper.}

\smallskip
\noindent \textit{नोट: पाँच प्रश्नों के उत्तर दीजिये जिनमें प्रश्न 1 अनिवार्य है। प्रश्न 1 के प्रत्येक भाग का उत्तर लगभग 200 शब्दों में तथा शेष प्रश्नों के उत्तर लगभग 1000 शब्दों में दीजिए। उपांत के अंक पूर्णांक के द्योतक हैं। अपने प्रश्नों को उचित मानचित्र एवं आरेख द्वारा प्रदर्शित कीजिए।}

\bigskip

\begin{enumerate}[label=\textbf{\arabic*.},leftmargin=*,itemsep=14pt]

  \item Write short notes on the following:\pts{4 $\times$ 3.5 = 14}
  
  निम्नलिखित पर संक्षिप्त टिप्पणियाँ लिखिए --
  \begin{parts}
    \item Concept of environmental ethics.\\
    पर्यावरणीय नैतिकता की अवधारणा।
    \item Concept of Deep Ecology.\\
    गहन पारिस्थितिकी
    \item Sustainable Development Goals - 13\\
    सतत् विकास लक्ष्य-13
    \item Ecological Niche.\\
    पारिस्थितिक निकेत
  \end{parts}

  \item Examine the Moral Theories of Environmental Ethics and explain its applicability in today's world.\pts{14}
  
  पर्यावरणीय नैतिकता के नैतिक सिद्धांतों का परीक्षण कीजिए तथा समकालीन विश्व में उनकी उपयुक्तता की व्याख्या कीजिए।

  \item Explain in detail the provisions of environmental legislations and laws in India, for natural resources protection.\pts{14}
  
  भारत में प्राकृतिक संसाधनों के संरक्षण हेतु निर्मित पर्यावरणीय विधानों एवं कानूनों के प्रावधानों का विस्तृत वर्णन कीजिए।

  \item Describe the major characteristics features of Indian Traditional Knowledge System and ecological values.\pts{14}
  
  भारतीय पारंपरिक ज्ञान प्रणाली की प्रमुख विशेषताओं तथा उसके पारिस्थितिक मूल्यों का वर्णन कीजिए।

  \item Explore the ethical, aesthetic and cultural values of environmental sustainability.\pts{14}
  
  पर्यावरणीय सततता के नैतिक, सौंदर्यात्मक तथा सांस्कृतिक मूल्यों का विश्लेषण कीजिए।

  \item Discuss in detail the major causes and impacts of habitat loss on biodiversity and the inherent challenges.\pts{14}
  
  जैव विविधता पर आपासीय क्षेत्र (Habitat) के क्षरण के प्रमुख कारणों एवं प्रभावों तथा उससे संबंधित अंतर्निहित चुनौतियों का विस्तारपूर्वक विवेचना कीजिए।

  \item Critically analyse how Geopolitics of food production shape the environmental sustainability and pursuit of Sustainable Development Goals.\pts{14}
  
  समालोचनात्मक विश्लेषण कीजिए कि खाद्य उत्पादन की भू-राजनीति किस प्रकार पर्यावरणीय सततता तथा सतत विकास लक्ष्यों की प्राप्ति को प्रभावित करती है।

\end{enumerate}

\bigskip
\begin{center}
  \textbf{*****}
\end{center}

\end{document}
